% Modernised reading — fonds Alexandre Grothendieck, shelfmark 42 (pages 1-6).
%
% This is an INTERPRETATION, not a transcription. It re-reads the pages in
% today's notation and vocabulary, states the hypotheses the manuscript leaves
% implicit, and drops the critical apparatus so that the mathematics reads
% straight through. Everything the brackets and underlines carried — a doubtful
% reading, a notation he used differently, a missing passage — moves into a
% footnote.
%
% It is held to being mathematically correct as it stands, not merely faithful
% to the page. Where the manuscript is loose or mistaken, the true statement is
% what appears here, and the footnote says what the page has. In any dispute of
% substance the transcription governs: whoever cites Grothendieck must cite
% batch-01.fr.tex.
%
% Scope: a SINGLE document for the whole folder, the convention every reading
% in this repository follows. The shelfmark holds one batch, so pages 1-6 are
% the folder entire and no part of it is read in another file. The folder was
% read whole — all six pages of batch-01.fr.tex — before any of this was
% written.
%
% Three departures from the page are worth naming here, because they are what
% a reader would otherwise trip on.
%
% 1. The folder orders its points twice, in opposite directions. Page 1 orders
%    X so that x <= y when y is a generization of x; page 3 orders it so that
%    y <= x when y is a generization of x. Both name the same relation. This
%    reading fixes one convention — specialization order, generizations
%    upward — and states the limit concretely, as coherent families, so that
%    nothing depends on the choice.
%
% 2. Lemme 3, and with it Corollaire 2 and the theorem of page 2, need Y to be
%    SEPARATED over S. The page never says so, and the argument it runs (two
%    morphisms agreeing on a schematically dense subscheme agree) is exactly
%    the step that fails without it. The affine line with a doubled origin is a
%    counterexample to Lemme 3 as stated; it is given below. Whether the
%    theorem itself survives without separatedness is not settled by these
%    pages, and this reading does not claim it either way.
%
% 3. The page glosses ker(A -> A_p) as "l'annulateur de A - p". That is not
%    what the kernel is; the correct description is given, with the page's own
%    in a footnote.
%
% Pass: Opus 5 (claude-opus-5), 2026-08-23 — first pass, unchecked against the
% transcription by a human.
\documentclass[11pt,a4paper]{article}
% Le préambule commun à toutes les transcriptions du fonds Grothendieck.
%
% Il tient en une idée : **l'incertitude se compose, elle ne se résout pas.**
% Une transcription qui lisse un mot douteux a détruit la seule chose pour
% laquelle on la faisait — un lecteur doit pouvoir savoir, à chaque endroit,
% ce qui était sur le feuillet et ce qui a été ajouté. Les six macros qui
% suivent sont donc l'appareil critique, et non de la décoration.
%
% Les mêmes commandes sont reconnues par `scripts/render.mjs`, qui produit la
% vue de lecture du site. Une macro ajoutée ici doit l'être aussi là-bas,
% faute de quoi le rendu échouera bruyamment — ce qui est voulu : un rendu
% silencieusement faux est pire qu'un rendu absent.

\ProvidesPackage{grothendieck}[2026/08/08 Transcriptions du fonds Grothendieck]

\usepackage{amsmath,amssymb,amsthm}
\usepackage{mathrsfs}
\usepackage{stmaryrd}
% Les diagrammes commutatifs sont partout dans ce fonds : c'est la langue
% même de la géométrie algébrique de Grothendieck, et une transcription qui
% les rendrait en prose ne transcrirait rien.
\usepackage{tikz-cd}
% Français par défaut — c'est la langue des feuillets — mais l'anglais est
% chargé aussi : la traduction et le résumé sont des documents anglais, et la
% typographie française (espaces avant les deux-points) y serait une faute.
% Ces documents basculent par \selectlanguage{english} en tête de corps.
\usepackage[english,french]{babel}
\usepackage{fontspec}
\usepackage{geometry}
\geometry{a4paper,margin=25mm,marginparwidth=28mm}
\usepackage{marginnote}
\usepackage{xcolor}
% ulem, not soul. `soul`'s \st and \ul are built on a character-by-character
% scan that XeLaTeX's font machinery breaks: the document fails with an
% undefined control sequence rather than degrading. ulem does the same two jobs
% and is Unicode-engine safe. `normalem` stops it hijacking \emph, which the
% transcriptions use for Grothendieck's own emphasis.
\usepackage[normalem]{ulem}
\usepackage{hyperref}
\hypersetup{colorlinks=true,linkcolor=black,urlcolor=black,pdfborder={0 0 0}}

% --- Métadonnées du lot -----------------------------------------------------
% Reprises telles quelles par le script de rendu, qui les affiche en tête de
% la vue de lecture. Elles fixent ce que le fichier prétend couvrir : sans
% elles, une transcription flotte sans point d'attache dans un fonds de seize
% mille feuillets.

\newcommand{\folder}[1]{\gdef\grothFolder{#1}}
\newcommand{\batch}[1]{\gdef\grothBatch{#1}}
\newcommand{\leaves}[2]{\gdef\grothFirst{#1}\gdef\grothLast{#2}}
\newcommand{\foldertitle}[1]{\gdef\grothFoldertitle{#1}}
\gdef\grothFolder{?}\gdef\grothBatch{?}\gdef\grothFirst{?}\gdef\grothLast{?}\gdef\grothFoldertitle{}

% --- Le résumé --------------------------------------------------------------
%
% Il ouvre la lecture modernisée, sous le titre « Résumé », et il est composé
% en petit. Ce n'est pas un ornement : c'est le seul passage du document qui
% ne soit pas des mathématiques, et un lecteur doit pouvoir le reconnaître
% avant de l'avoir lu — puis le sauter s'il connaît déjà le sujet, ou s'y
% arrêter s'il ne le connaît pas. Le filet qui le suit marque la reprise du
% corps.

\newenvironment{resume}
  {\small\setlength{\parskip}{0.45em}}
  {\par\medskip\hrule\medskip}

% --- Mots-clés ---------------------------------------------------------------
%
% En anglais, en clôture du résumé de la lecture modernisée : le vocabulaire
% moderne sous lequel chercher ce que les feuillets construisent. Le manifeste
% du site (`npm run manifest`) les extrait pour étiqueter la cote entière —
% cette ligne est la source unique des tags, il n'y a pas de fichier à côté.
\newcommand{\keywords}[1]{\par\smallskip\noindent
  {\footnotesize\textsc{Keywords} --- \textit{#1}}\par}

% --- L'appareil critique ----------------------------------------------------

% Le numéro de feuillet, en marge. C'est l'ancre qui permet de régler un
% désaccord sur une formule en regardant *un* feuillet plutôt qu'une chemise
% de six cents. Le site s'en sert aussi pour tourner les pages du fac-similé
% quand on fait défiler la transcription.
\newcommand{\leaf}[1]{%
  \marginnote{\footnotesize\color{gray!70}\textsf{#1}}%
  \label{leaf:#1}\ignorespaces}

% Illisible : on ne devine pas. Un mot inventé qui se lit comme les autres est
% le pire résultat possible.
\newcommand{\ill}{\textcolor{red!60!black}{[\,\ldots\,]}}

% Lecture douteuse : le mot est donné, et signalé comme incertain.
\newcommand{\uncertain}[1]{\uline{#1}}

% Ajout de l'éditeur — une lettre manquante, un mot sous-entendu.
\newcommand{\add}[1]{\textcolor{blue!50!black}{[#1]}}

% Biffé par Grothendieck lui-même. Ce qu'il a rayé fait partie du document :
% une rature dit souvent où la pensée a changé de direction.
\newcommand{\struck}[1]{\sout{#1}}

% Note du transcripteur, sur l'état matériel du feuillet.
\newcommand{\note}[1]{\par\smallskip{\small\color{gray!60!black}\textit{[#1]}}\par\smallskip}

% Note marginale de Grothendieck, distincte du corps du texte.
\newcommand{\marginal}[1]{\par\smallskip\noindent\hspace{1em}%
  {\small\color{gray!60!black}#1}\par\smallskip}

% --- Titre ------------------------------------------------------------------

\AtBeginDocument{%
  \begin{center}
    {\large\bfseries Fonds Alexandre Grothendieck --- cote n\textsuperscript{o}~\grothFolder}\\[2pt]
    {\itshape\grothFoldertitle}\\[4pt]
    {\small feuillets \grothFirst--\grothLast\ (lot \grothBatch)}\\[2pt]
    {\footnotesize\color{gray!60!black}Transcription établie d'après le fac-similé numérisé
     par l'Université de Montpellier.\\
     Les lectures incertaines sont soulignées, les illisibles notées [\,\ldots\,],
     les ajouts entre crochets.}
  \end{center}
  \vspace{1em}}

\endinput


\folder{42}
\pages{1}{6}
\dating{[années 1960-1970]}
\watermark{Édition de démonstration\\interprétation personnelle de l'œuvre}
\foldertitle{Réseaux d'anneaux locaux : notes manuscrites (s.d.) — lecture modernisée du dossier entier}

\begin{document}

\section*{Résumé}

\begin{resume}
Un schéma --- Grothendieck écrit \emph{préschéma} --- se décrit de deux
façons. On peut le regarder comme un espace : des ouverts, et sur chaque
ouvert un anneau de fonctions. Ou bien on peut ne retenir que trois choses
beaucoup plus pauvres : la liste de ses points ; la relation « le point $y$
est adhérent au point $x$ », qui dit lequel est plus spécial et lequel plus
générique ; et, pour chaque point, l'anneau des fonctions qu'on voit depuis ce
point d'aussi près qu'on veut. Ces six pages demandent si cette seconde
description, qui a jeté la topologie, suffit encore. Grothendieck appelle
\emph{réseau d'anneaux locaux} la donnée appauvrie, et démontre que oui : une
application entre deux schémas est exactement une application entre leurs
réseaux.

Le problème qui se pose est le suivant, et il est plus subtil qu'il n'en a
l'air. Qu'une application soit \emph{déterminée} par ce qu'elle fait au
voisinage de chaque point, c'est banal : deux fonctions qui coïncident près de
chaque point coïncident. Qu'on puisse en sens inverse \emph{fabriquer} une
application en la donnant point par point, c'est autre chose. On se donne, en
chaque point, un bout d'application défini tout autour de lui, et l'on exige
seulement que deux bouts voisins se raccordent au point plus générique des
deux. Rien ne dit que ces bouts se recollent en une seule application : il
faut que le lieu où deux d'entre eux coïncident soit ouvert, c'est-à-dire ait
de l'épaisseur, et non pas réduit à quelques points isolés.

L'idée qui règle cela est le c\oe ur du dossier, et elle est jolie. Le lieu où
un bout d'application coïncide avec ce qui était prescrit a toujours une
propriété de stabilité gratuite : s'il contient un point, il contient tous les
points plus génériques que lui. Cette stabilité seule ne fait pas un ouvert.
Mais si on lui adjoint une seconde propriété --- être \emph{constructible},
c'est-à-dire s'obtenir en un nombre fini d'opérations élémentaires sur des
ouverts et des fermés --- alors le lieu est ouvert, et tout se recolle. Reste
à obtenir la constructibilité, et c'est un fait d'algèbre commutative qui la
donne : dans un anneau noethérien $A$, le passage au localisé $A_{\mathfrak
p}$ perd de l'information, mais on peut inverser \emph{un seul} élément
$f$ hors de $\mathfrak p$ après quoi il n'en perd plus. Un seul élément :
c'est ce « un seul » qui fabrique l'ouvert.

Ces pages montrent aussi comment il travaillait. La page 2 porte une première
démonstration, menée directement sur les réseaux, entièrement barrée de deux
grandes croix ; la page 3 recommence à un tout autre étage, celui des
faisceaux et de l'algèbre commutative, où le contenu réel se trouvait. Rien
n'est effacé, la tentative abandonnée reste sur la feuille. Et la dernière
page se termine par un contre-exemple, c'est-à-dire par la vérification que
l'hypothèse coûteuse qu'il a dû introduire ne pouvait pas être évitée.

Le vocabulaire d'aujourd'hui pour ce cercle d'idées : l'ordre de
spécialisation d'un espace sobre, les parties constructibles et le critère
« constructible et stable par générisation, donc ouvert », les limites de
fibres le long de cet ordre, et --- pour la partie que le manuscrit laisse
implicite --- les morphismes schématiquement dominants et la séparation.

\keywords{specialization order, sober space, constructible set, limit of stalks, locally ringed space, schematically dominant morphism}
\end{resume}

\pagerange{1}{2}
\section*{Le réseau d'un préschéma, et l'énoncé (pages 1 et 2)}

Soit $X$ un espace topologique. Pour $x, y \in X$ on écrit $y \rightsquigarrow
x$, et l'on dit que $y$ est une \emph{spécialisation} de $x$, ou que $x$ est
une \emph{générisation} de $y$, lorsque $y \in \overline{\{x\}}$. C'est une
relation de préordre, et un ordre dès que $X$ est $T_0$.

\begin{quote}
\textbf{Définition.} Un \emph{réseau d'anneaux locaux} est la donnée d'un
ensemble ordonné $X$ et, pour chaque $x \in X$, d'un anneau local $O_x$,
munis pour chaque couple $y \rightsquigarrow x$ d'un homomorphisme
$O_y \to O_x$, fonctoriel, tel que
\[ (O_y)_{\mathfrak p_{yx}} \;\xrightarrow{\ \sim\ }\; O_x, \]
où $\mathfrak p_{yx} \subset O_y$ est l'image réciproque de l'idéal maximal
de $O_x$.\footnote{Le mot que nous lisons « réciproque » est illisible sur la
page 1, qui porte seulement « $p_{xy}$ est l'image \dots\ de
$\mathfrak m(O_y)$ ». C'est la seule lecture qui donne un idéal premier de la
source, donc la seule sous laquelle la condition ait un sens ; la
reconstruction est la nôtre.} Autrement dit : chaque anneau plus générique est un localisé de
chaque anneau plus spécial, en le premier idéal qu'il désigne.
\end{quote}

\noindent
Un \emph{homomorphisme de réseaux} $X \to Y$ est une application croissante
$f : X \to Y$ accompagnée, pour chaque $x$, d'un homomorphisme \emph{local}
$O_{f(x)} \to O_x$, le tout compatible aux applications de
générisation.\footnote{Page 1. Le manuscrit hésite sur ce dernier mot : il
écrit d'abord « spécialisation », le biffe, et met « généralisation ». Le sens
est celui-ci : la compatibilité porte sur les carrés indexés par
$y \rightsquigarrow x$.}

La condition de localisation a une conséquence que Grothendieck note en marge
de la page 1 : pour $x$ fixé, les générisations de $x$ correspondent
bijectivement aux idéaux premiers de $O_x$. C'est la raison d'être de la
définition --- le réseau retient assez de l'espace pour que le spectre de
chaque anneau local se lise dans l'ordre lui-même.

Tout préschéma donne un réseau : ses points, leur ordre de spécialisation, ses
anneaux locaux. Cela définit un foncteur
\[ p \,:\, \mathbf{Sch} \longrightarrow \mathbf{Res}. \]

\begin{quote}
\textbf{Théorème (page 2).} Soit $S$ un préschéma localement noethérien. La
restriction de $p$ aux $S$-préschémas \emph{localement de type fini} est
\emph{pleinement fidèle} à valeurs dans les réseaux sur $p(S)$ :
\[ \mathrm{Hom}_S(X, Y) \;\xrightarrow{\ \sim\ }\;
   \mathrm{Hom}_{p(S)}\bigl(p(X), p(Y)\bigr). \]
\end{quote}

\noindent
C'est ce que le dossier entier démontre.\footnote{Sous une hypothèse de
séparation que la page ne formule pas ; voir la section « Des morphismes, et
non plus des fonctions » et la note qui s'y trouve. L'énoncé est reproduit
ici tel que la page 2 le donne, et corrigé là où l'argument l'exige.} La
page 2 en contient aussi une première démonstration, tentée directement sur
les réseaux et barrée de deux grandes croix : elle passait par un dévissage
des points de $X$ par spécialisation et un ensemble fini de recouvrements,
puis par la construction d'un homomorphisme d'un voisinage de $x$ dans un
voisinage de $f(x)$. Elle est abandonnée en cours de route, et la
démonstration reprend à neuf page 3, à l'étage des faisceaux.

\medskip
\noindent
\textbf{Une convention, fixée ici pour tout ce qui suit.} Le manuscrit ordonne
$X$ deux fois, dans deux sens opposés : la page 1 pose $x \leq y$ lorsque $y$
est une générisation de $x$, la page 3 pose $y \leq x$ dans le même cas. Les
deux nomment la même relation ; seule la flèche du symbole $\leq$ change de
sens.\footnote{Page 1 : « un système inductif d'anneaux locaux
$(O_x)$ \dots tel que pour $x \leq y$, $O_x \to O_y$ induise un isomorphisme
$(O_x)_{p_{xy}} \to O_y$ », et la marge « une correspondance bijective entre
les $y \geq x$ et les idéaux premiers de $O_x$ » --- donc générique vers le
haut. Page 3 : « Ordonnons $X$ par $y \leq x \Leftrightarrow \bar y \ni x$ »
--- donc générique vers le bas. C'est une inadvertance, non une seconde
notion. Cette lecture n'emploie ni l'un ni l'autre symbole et écrit partout
$y \rightsquigarrow x$, puis décrit les limites comme des familles cohérentes,
ce qui est indépendant du choix.}

\pagerange{3}{3}
\section*{Une section globale est une famille cohérente de germes (page 3)}

Soit $X$ un espace topologique et $F$ un faisceau d'ensembles sur $X$. Pour
$y \rightsquigarrow x$, tout ouvert contenant $y$ contient $x$, de sorte qu'il
y a une application canonique entre fibres
\[ \rho_{yx} \,:\, F_y \longrightarrow F_x, \]
du point spécial vers le point générique. Appelons \emph{famille cohérente}
une famille $(\xi_x)_{x \in X}$ avec $\xi_x \in F_x$ telle que
$\rho_{yx}(\xi_y) = \xi_x$ chaque fois que $y \rightsquigarrow x$. Les
familles cohérentes forment la limite $\varprojlim_{x} F_x$ du diagramme des
fibres indexé par l'ordre de spécialisation, et il y a une application
canonique $\Gamma(X, F) \to \varprojlim_x F_x$, qui envoie une section sur la
famille de ses germes.

\begin{quote}
\textbf{Proposition.} Soit $X$ un espace localement noethérien et
\emph{sobre} --- toute partie fermée irréductible a exactement un point
générique. Soit $F$ un faisceau sur $X$ vérifiant :
\begin{quote}
$(\mathrm{H})$ \quad pour tout $x \in X$, il existe un ouvert non vide $U$ de
$\overline{\{x\}}$ tel que $\rho_{yx} : F_y \to F_x$ soit \emph{injective}
pour tout $y \in U$.
\end{quote}
Alors $\Gamma(X, F) \to \varprojlim_x F_x$ est bijective.
\end{quote}

\noindent
L'injectivité ne demande rien : une section dont tous les germes sont égaux à
ceux d'une autre leur est égale, par l'axiome de séparation des faisceaux.
Tout le contenu est dans la surjectivité, et il tient dans le lemme
suivant.\footnote{Sur le mot « injective », qui gouverne $(\mathrm{H})$ :
il est écrit d'un seul trait sur la page, la première syllabe réduite à un
crochet, et se lirait aussi bien « surjective ». C'est l'usage qui tranche ---
le lemme 2 de la page 4 fournit une injectivité, et la page 5 en conclut que
« $F_{\mathfrak q} \to F_{\mathfrak p}$ est injectif » vérifie la condition de
la proposition. La transcription porte la même note à la première occurrence.}

\begin{quote}
\textbf{Lemme 1.} Soit $(\xi_x)$ une famille cohérente et $s \in F(U)$ une
section sur un ouvert $U$. Alors
$E = \{\, x \in U \;:\; s_x = \xi_x \,\}$
est \emph{ouvert}.
\end{quote}

\noindent
La démonstration se fait en deux temps, et c'est leur conjonction qui est
l'idée du dossier.

\emph{$E$ est stable par générisation.} Soit $x \in E$ et soit $x'$ une
générisation de $x$ dans $U$. Alors $x \rightsquigarrow x'$, l'application
$\rho_{x x'} : F_x \to F_{x'}$ est définie, et elle envoie $s_x$ sur $s_{x'}$
et $\xi_x$ sur $\xi_{x'}$. De $s_x = \xi_x$ suit donc $s_{x'} = \xi_{x'}$,
c'est-à-dire $x' \in E$.

\emph{$E$ est constructible.} On peut supposer $X$ noethérien, la question
étant locale. Soit $Z \subseteq U$ une partie fermée irréductible, de point
générique $z$. Ou bien $z \notin E$, et alors $E \cap Z = \varnothing$ : car
tout point de $Z$ a $z$ pour générisation, et $E$ étant stable par
générisation, un point de $E \cap Z$ forcerait $z \in E$. Ou bien $z \in E$,
et alors $(\mathrm{H})$ fournit un ouvert non vide $V$ de $\overline{\{z\}} =
Z$ sur lequel $\rho_{yz} : F_y \to F_z$ est injective ; pour $y \in V$ on a
$\rho_{yz}(s_y) = s_z = \xi_z = \rho_{yz}(\xi_y)$, d'où $s_y = \xi_y$ par
injectivité, donc $V \subseteq E \cap Z$.\footnote{C'est le pas où
l'hypothèse $(\mathrm{H})$ paie, et le seul endroit du dossier où elle serve.
La page écrit ici « d'où $s_x = \xi_y$ » ; l'égalité dont l'argument a besoin,
et qu'il obtient, est $s_y = \xi_y$. On a tenu la première pour un lapsus de
plume.} Le critère de constructibilité dans un espace noethérien est
satisfait dans les deux cas.

Une partie constructible et stable par générisation d'un espace noethérien est
ouverte.\footnote{Résultat standard, contemporain de ces notes ; il figure
dans EGA IV, au chapitre des parties constructibles. Le manuscrit l'emploie
sans le nommer --- « il suffit de montrer qu'il est constructible » --- ce qui
est le signe qu'il l'avait sous la main.} Donc $E$ est ouvert, et le lemme 1
est acquis.

La proposition suit. Soit $(\xi_x)$ une famille cohérente et $x_0 \in X$. Le
germe $\xi_{x_0}$ provient d'une section $s$ sur un ouvert $U \ni x_0$ ; par
le lemme 1, $E = \{x \in U : s_x = \xi_x\}$ est ouvert et contient $x_0$. La
famille est donc réalisée par une section au voisinage de chaque point, et
deux telles sections ont, sur leur intersection, les mêmes germes en tout
point : elles y coïncident. Elles se recollent en une section globale dont la
famille de germes est $(\xi_x)$.

\pagerange{4}{5}
\section*{Pourquoi le faisceau structural vérifie $(\mathrm{H})$ (pages 4 et 5)}

L'application à retenir est celle du faisceau structural, et elle exige un
fait d'algèbre commutative que le dossier isole proprement.

\begin{quote}
\textbf{Lemme 2.} Soient $A$ un anneau noethérien et $\mathfrak p$ un idéal
premier de $A$. Alors :
\begin{enumerate}
\item il existe $f \in A - \mathfrak p$ tel que $A_f \to A_{\mathfrak p}$ soit
      injectif ;
\item si de plus $A \to A_{\mathfrak p}$ est déjà injectif, alors
      $A_{\mathfrak q} \to A_{\mathfrak p}$ est injectif pour tout idéal
      premier $\mathfrak q \supseteq \mathfrak p$.
\end{enumerate}
\end{quote}

\noindent
\emph{Démonstration de (1).} Soit $J = \ker(A \to A_{\mathfrak p})$,
c'est-à-dire l'ensemble des $a \in A$ annulés par \emph{au moins un} élément
de $A - \mathfrak p$.\footnote{La page glose $J$ comme « l'annulateur de
$A - \mathfrak p$ ». Ce n'est pas ce qu'est le noyau : l'annulateur de
$A - \mathfrak p$ serait l'ensemble des $a$ annulés par \emph{tout} élément de
$A - \mathfrak p$, ensemble en général bien plus petit et qui ne conviendrait
pas ici. La démonstration qui suit sur la page utilise, elle, la bonne
description --- elle prend un annulateur par générateur.} L'anneau étant
noethérien, $J = (x_1, \dots, x_n)$, et chaque $x_i$ est annulé par un
$s_i \in A - \mathfrak p$. Posons $f = s_1 \cdots s_n$ : il est hors de
$\mathfrak p$, qui est premier, et $fJ = 0$. Soit alors $x/f^n$ dans le noyau
de $A_f \to A_{\mathfrak p}$ : il existe $s \in A - \mathfrak p$ avec
$sx = 0$, donc $x \in J$, donc $fx = 0$, donc $x/f^n = 0$ dans $A_f$.

\emph{Démonstration de (2).} Comme $\mathfrak p \subseteq \mathfrak q$, tout
élément de $A - \mathfrak q$ est hors de $\mathfrak p$, et
$A_{\mathfrak q} \to A_{\mathfrak p}$ a un sens. Soit $x/s$ ($x \in A$,
$s \in A - \mathfrak q$) dans son noyau : l'image de $x$ dans $A_{\mathfrak p}$
est nulle, donc $x = 0$ par l'injectivité supposée de $A \to A_{\mathfrak p}$,
donc $x/s = 0$.\footnote{La page traite ce point sous la forme : « on peut
supposer $f = 1$, i.e.\ $A_f = A$, car $A_{\mathfrak p} =
(A_f)_{\mathfrak p A_f}$ ». C'est le même argument, la réduction étant
légitime parce que l'injectivité de $A \to A_{\mathfrak p}$ entraîne celle de
$A_h \to A_{\mathfrak p}$ pour tout $h$.}

\begin{quote}
\textbf{Corollaire 1.} Soit $X$ un préschéma localement noethérien. Alors
\[ \Gamma(X, \mathcal{O}_X) \;\xrightarrow{\ \sim\ }\;
   \varprojlim_{x \in X} \mathcal{O}_{X,x}, \]
c'est-à-dire : une fonction globale sur $X$ est exactement la donnée d'une
famille $(\xi_x) \in \prod_x \mathcal{O}_{X,x}$ telle que $\xi_x$ soit
l'image de $\xi_y$ par la générisation chaque fois que
$y \rightsquigarrow x$.
\end{quote}

\noindent
Il suffit de vérifier $(\mathrm{H})$ pour $\mathcal{O}_X$, et c'est exactement
le lemme 2. La question étant locale, prenons $X = \operatorname{Spec} A$ avec
$A$ noethérien, et $x = \mathfrak p$. Alors $\overline{\{x\}} = V(\mathfrak p)$,
et le lemme 2 (1) donne $f \notin \mathfrak p$ tel que
$A_f \to A_{\mathfrak p}$ soit injectif ; le lemme 2 (2), appliqué à $A_f$,
donne l'injectivité de $A_{\mathfrak q} \to A_{\mathfrak p}$ pour tout
$\mathfrak q \in D(f) \cap V(\mathfrak p)$. Or $D(f) \cap V(\mathfrak p)$ est
un ouvert non vide de $\overline{\{x\}}$ --- il contient $\mathfrak p$
lui-même. C'est $(\mathrm{H})$.

\pagerange{5}{6}
\section*{Des morphismes, et non plus des fonctions (pages 5 et 6)}

Le pas suivant remplace le faisceau structural par un faisceau de morphismes,
et c'est lui qui donne le théorème.

Soient $S$ un préschéma localement noethérien, $X$ un $S$-préschéma
localement noethérien et $Y$ un $S$-préschéma. On pose
\[ F(U) \;=\; \mathrm{Hom}_S(U, Y) \qquad (U \subseteq X \text{ ouvert}), \]
qui est un faisceau puisque les morphismes se recollent, et dont les sections
globales sont $\mathrm{Hom}_S(X, Y)$ par construction.

\begin{quote}
\textbf{Lemme 3.} Supposons $Y$ \emph{séparé} sur $S$. Alors $F$ vérifie
$(\mathrm{H})$. Plus précisément, si $X = \operatorname{Spec} A$ et si
$A \to A_{\mathfrak p}$ est injectif, alors $F_{\mathfrak q} \to
F_{\mathfrak p}$ est injective pour tout $\mathfrak q \supseteq \mathfrak p$.
\end{quote}

\noindent
\emph{Démonstration.} Soient deux germes en $\mathfrak q$, représentés par
$f, g : U \to Y$ avec $U \ni \mathfrak q$ ouvert, de mêmes germes en
$\mathfrak p$.\footnote{Noter que $\mathfrak p \in U$ : un ouvert est stable
par générisation, et $\mathfrak p \subseteq \mathfrak q$ fait de $\mathfrak p$
une générisation de $\mathfrak q$. La page écrit ici « pour les $\mathfrak q$
générisant $\mathfrak p$ » ; la condition dont l'argument a besoin, et la
seule sous laquelle $A_{\mathfrak q} \to A_{\mathfrak p}$ existe, est
$\mathfrak q \supseteq \mathfrak p$, c'est-à-dire que $\mathfrak q$
\emph{spécialise} $\mathfrak p$. Le sens est inversé sur la page.} Quitte à
rétrécir $U$ en un ouvert principal, on peut supposer $U = \operatorname{Spec}
A$ avec $A \to A_{\mathfrak p}$ injectif. Il existe un ouvert $W \ni
\mathfrak p$ sur lequel $f$ et $g$ coïncident. L'égalisateur $E$ de $f$ et $g$
est le produit fibré de $U \to Y \times_S Y$ le long de la diagonale : c'est
une immersion \emph{fermée} précisément parce que $Y$ est séparé sur $S$. Or
$\operatorname{Spec} A_{\mathfrak p} \to \operatorname{Spec} A$ se factorise
par $W$, donc par $E$ ; et cette flèche est schématiquement dominante, puisque
$A \to A_{\mathfrak p}$ est injectif. Un sous-schéma fermé par lequel se
factorise un morphisme schématiquement dominant est le schéma tout entier :
$E = U$, donc $f = g$. \hfill$\square$

\begin{quote}
\textbf{Corollaire 2.} Sous les mêmes hypothèses, la donnée d'un
$S$-morphisme $X \to Y$ équivaut à celle d'une famille cohérente de germes
$f_x \in F_x$, $x \in X$.
\end{quote}

\noindent
C'est la proposition appliquée à $F$. Une forme équivalente, que la page 6
donne séparément : si $f, g : \operatorname{Spec} A \to Y$ sont deux
$S$-morphismes dont les composés avec $\operatorname{Spec} A_{\mathfrak p} \to
\operatorname{Spec} A$ coïncident, et si $A \to A_{\mathfrak p}$ est injectif,
alors $f = g$.

\begin{quote}
\textbf{Corollaire 3.} Soient $S$ localement noethérien, $X$ et $Y$ des
$S$-préschémas localement de type fini, $Y$ séparé sur $S$. Alors
\[ \mathrm{Hom}_S(X, Y) \;\xrightarrow{\ \sim\ }\;
   \mathrm{Hom}_{p(S)}\bigl(p(X), p(Y)\bigr), \]
ce qui est le théorème de la page 2.
\end{quote}

\noindent
L'injectivité est le corollaire 2. Pour la surjectivité, un homomorphisme de
réseaux $p(X) \to p(Y)$ fournit, en chaque $x$, un point $\varphi(x) \in Y$ et
un homomorphisme local $\mathcal{O}_{Y, \varphi(x)} \to \mathcal{O}_{X,x}$,
c'est-à-dire un $S$-morphisme $\operatorname{Spec} \mathcal{O}_{X,x} \to Y$.
Pour en faire un germe --- un morphisme défini sur tout un voisinage de $x$
--- on utilise que
\[ \varinjlim_{U \ni x} \mathrm{Hom}_S(U, Y)
   \;\xrightarrow{\ \sim\ }\;
   \mathrm{Hom}_S\bigl(\operatorname{Spec} \mathcal{O}_{X,x},\, Y\bigr), \]
ce qui vaut dès que $Y$ est localement de présentation finie sur $S$ ---
c'est-à-dire ici, $S$ étant localement noethérien, localement de type
fini.\footnote{Théorèmes de passage à la limite d'EGA IV. C'est le seul
endroit où l'hypothèse « de type fini » sur $Y$ sert, et c'est pourquoi la
page 6 la fait apparaître au corollaire 3 et non au corollaire 2 : le
corollaire 2 n'en a pas besoin. Grothendieck note d'ailleurs en marge de son
corollaire 3, à propos de l'hypothèse noethérienne sur $S$, « ce dernier
est-il nécessaire ? » --- question qu'il laisse ouverte.} Les germes ainsi
obtenus forment une famille cohérente, et le corollaire 2 les recolle.

\medskip
\noindent
\textbf{La séparation, que la page ne dit pas.} Aucun des énoncés du
manuscrit ne suppose $Y$ séparé sur $S$, et le lemme 3 est faux sans cette
hypothèse. Prenons $S = \operatorname{Spec} k$, $X = \mathbb{A}^1_k =
\operatorname{Spec} k[T]$, et pour $Y$ la droite affine à origine dédoublée.
Les deux inclusions $j_1, j_2 : \mathbb{A}^1 \to Y$ coïncident hors de
l'origine. Au point générique $\mathfrak p = (0)$ elles ont donc le même
germe ; à l'origine $\mathfrak q = (T)$ elles ont des germes distincts,
puisqu'elles y envoient le point sur deux points différents de $Y$ et
diffèrent donc sur tout voisinage. L'application
$F_{\mathfrak q} \to F_{\mathfrak p}$ n'est pas injective, alors même que
$X$ est intègre et que $k[T] \to k(T)$ est injectif : la mise en défaut vient
tout entière de $Y$.\footnote{Ce contre-exemple est le nôtre ; le manuscrit ne
soulève pas la question. Il montre que le lemme 3 tel qu'il est énoncé est
faux, non que le théorème le soit : la surjectivité du corollaire 2 peut
survivre à la perte du lemme 1, et l'exemple ci-dessus n'en fournit pas de
contre-exemple. Ce que ces six pages établissent, on l'a donc énoncé avec
l'hypothèse de séparation ; ce qu'il en resterait sans elle, elles ne le
disent pas, et cette lecture ne le tranche pas.}

\pagerange{6}{6}
\section*{La condition ne peut pas être supprimée (page 6)}

Le dossier se termine par la vérification que $(\mathrm{H})$ n'est pas une
commodité de démonstration. Sans elle, la proposition tombe --- et elle tombe
du côté de la surjectivité, c'est-à-dire là où était tout le contenu.

Soit $X = \operatorname{Spec} k[T] = \mathbb{A}^1_k$, avec $k$ algébriquement
clos, de sorte que les points fermés sont les $a \in k$ et qu'il y a un unique
point générique $\eta$. Pour chaque point fermé $a$ soit $E_a$ un ensemble
pointé non trivial, et soit
\[ F \;=\; \bigoplus_{a} \, i_{a *} E_a \]
la somme directe des faisceaux gratte-ciel correspondants.\footnote{La page
écrit $F = \coprod_{x \in X} E_x$, sans préciser si $x$ parcourt tous les
points ou seulement les points fermés ; c'est la seconde lecture qui fait
marcher l'exemple, et seuls les points fermés y contribuent de toute façon.
La fin de la page 6 est très fragmentaire dans le fac-similé --- une ligne sur
deux se laisse lire --- et la reconstruction qui suit est la nôtre, faite le
long de ce qui reste : « l'élément \emph{unité} de $\Gamma(E_x)$ », «
$\xi_a = 0$ », « les $\xi_x$ ($x \neq x_i$) forment un système cohérent ».}

Les fibres sont $F_a = E_a$ en un point fermé et $F_\eta = 0$ au point
générique. La condition de cohérence est donc vide : \emph{toute} famille
$(\xi_a) \in \prod_a E_a$ est cohérente. Mais les sections globales d'une
somme directe de gratte-ciel sont les familles à support \emph{fini} --- sur
$\mathbb{A}^1$ tout ouvert non vide est cofini, et une famille localement
finie y est finie. Prenant $\xi_a$ égal à l'élément distingué de $E_a$ pour
tout $a$ sauf un, où on le prend nul, on obtient une famille cohérente à
support infini : elle ne provient d'aucune section globale.

L'hypothèse $(\mathrm{H})$ est bien en défaut : elle demanderait, pour
$x = \eta$, un ouvert non vide $U$ de $\overline{\{\eta\}} = X$ sur lequel
$F_a \to F_\eta = 0$ soit injective, ce qui forcerait $E_a = 0$.

\medskip
\noindent
C'est là que le dossier s'arrête. L'exemple ne dit pas seulement que la
démonstration avait besoin de $(\mathrm{H})$ : il dit que le passage du local
au global, pour un faisceau quelconque sur un schéma, n'est pas donné par
l'ordre de spécialisation seul. Ce que le corollaire 1 établit pour
$\mathcal{O}_X$, et le corollaire 3 pour les morphismes, tient à une propriété
de finitude de l'algèbre commutative --- un seul $f$ à inverser --- et non à
la forme de l'ordre.

\end{document}
